\subsection{PARABOLIC SUBALGEBRAS}

PROPOSITION 11. Let \(\mathfrak{p} = \mathfrak{h} + \mathfrak{g}^{\mathrm{P}}\) be a subalgebra of \(\mathfrak{g}\) containing \(\mathfrak{h}\) . The following conditions are equivalent:

\begin{enumerate}
\def\labelenumi{(\roman{enumi})}
\item
  p contains a Borel subalgebra of \((\mathfrak{g}, \mathfrak{h})\) ;
\item
  there exists a chamber C of R such that \(\mathrm{P} \supset \mathrm{R}_{+}(\mathrm{C})\) ;
\item
  P is parabolic, in other words (Chap. VI, §1, no. 7, Def. 4), \(\mathrm{P} \cup (-\mathrm{P}) = \mathrm{R}\) .
\end{enumerate}

Conditions (i) and (ii) are equivalent by Prop. 7. Conditions (ii) and (iii) are equivalent by Chap. VI, §1, no. 7, Prop. 20.

DEFINITION 2. A subalgebra of g containing h and satisfying the equivalent conditions of Prop. 11 is called a parabolic subalgebra of \((\mathfrak{g}, \mathfrak{h})\) . A parabolic subalgebra of g is a parabolic subalgebra of \((\mathfrak{g}, \mathfrak{h}')\) where \(h'\) is a splitting Cartan subalgebra of g.

This definition extends immediately to the case in which \((\mathfrak{g}, \mathfrak{h})\) is a split reductive Lie algebra.

Remark. Let B be a basis of R, and b the corresponding Borel subalgebra. If \(\Sigma \subset B\) , denote by \(Q_{\Sigma}\) the set of roots that are linear combinations of elements of \(\Sigma\) with coefficients \(\leq 0\) ; put \(\mathfrak{p}(\Sigma) = \mathrm{R}_{+}(\mathrm{B}) \cup \mathrm{Q}_{\Sigma}\) and \(\mathfrak{p}_{\Sigma} = \mathfrak{h} \oplus \mathfrak{g}^{\mathrm{P}(\Sigma)}\) . By Chap. VI, §1, no. 7, Lemma 3 and Prop. 20, \(p_{\Sigma}\) is a parabolic subalgebra containing b and every parabolic subalgebra of g containing b is obtained in this way.

Lemma 3. Let V be a finite dimensional real vector space, S a root system in \(V^{*}\) , P the set of parabolic subsets of S; let H be the set of Ker \(\alpha\) for \(\alpha \in S\) , and F the set of facets of V relative to H (Chap. V, §1, no. 2, Def. 1).

If \(\mathrm{P} \in \mathcal{P}\) , let \(\overline{\mathrm{F}}(\mathrm{P})\) be the set of \(v \in \mathrm{V}\) such that \(\alpha(v) \geq 0\) for all \(\alpha \in \mathrm{P}\) . If \(\mathrm{F} \in \mathcal{F}\) , let \(\mathrm{P}(\mathrm{F})\) be the set of \(\alpha \in \mathrm{R}\) such that \(\alpha(v) \geq 0\) for all \(v \in \mathrm{F}\) .

Then \(\mathrm{F} \mapsto \mathrm{P}(\mathrm{F})\) is a bijection from \(\mathcal{F}\) to \(\mathcal{P}\) ; for all \(\mathrm{F} \in \mathcal{F}\) , \(\overline{\mathrm{F}}(\mathrm{P}(\mathrm{F}))\) is the closure of \(\mathrm{F}\) .

\begin{enumerate}
\def\labelenumi{\alph{enumi})}
\tightlist
\item
  Let \(P \in P\) . There exists a chamber C of S and a subset \(\Sigma\) of the basis \(\mathrm{B}(\mathrm{C})\) such that \(\mathrm{P} = \mathrm{S}_{+}(\mathrm{C}) \cup \mathrm{Q}\) where Q is the set of linear combinations of elements of \(\Sigma\) with non-positive integer coefficients (Chap. VI, §1, no. 7, Prop. 20). Put
\end{enumerate}

\[
\mathrm{B} (\mathrm{C}) = \{\alpha_ {1}, \dots , \alpha_ {l} \}, \Sigma = \{\alpha_ {1}, \dots , \alpha_ {m} \}.
\]

If \(v \in \mathrm{V}\) , we have the following equivalences:

\[
\begin{array}{c} \alpha (v) \geq 0 \text {   for   all   } \alpha \in \mathrm{P} \\ \Longleftrightarrow \alpha_ {1} (v) \geq 0, \ldots \alpha_ {l} (v) \geq 0, \alpha_ {1} (v) \leq 0, \ldots , \alpha_ {m} (v) \leq 0 \\ \Longleftrightarrow \alpha_ {1} (v) = \dots = \alpha_ {m} (v) = 0, \alpha_ {m + 1} (v) \geq 0, \ldots , \alpha_ {l} (v) \geq 0, \end{array}
\]

so \(\overline{\mathrm{F}}(\mathrm{P})\) is the closure of the set

\[
\{v \in \mathrm{V} \mid \alpha_ {1} (v) = \dots = \alpha_ {m} (v) = 0, \alpha_ {m + 1} (v) > 0, \dots , \alpha_ {l} (v) > 0 \},
\]

a set which is a facet F relative to H since every element of S is a linear combination of \(\alpha_{1},\ldots,\alpha_{l}\) in which the coefficients are either all \(\geq0\) or all \(\leq0\) . Moreover, if \(\beta=u_{1}\alpha_{1}+\cdots+u_{l}\alpha_{l}\in S\) ,

\[
\begin{array}{c} \beta \in \mathrm{P} (\mathrm{F}) \Longleftrightarrow u _ {m + 1} \geq 0, \ldots , u _ {l} \geq 0 \\ \Longleftrightarrow \beta \in \mathrm{S} _ {+} (\mathrm{C}) \text {   or   } (- \beta \in \mathrm{S} _ {+} (\mathrm{C}) \text {   and   } u _ {m + 1} = \ldots = u _ {l} = 0) \\ \Longleftrightarrow \beta \in \mathrm{S} _ {+} (\mathrm{C}) \cup \mathrm{Q} = \mathrm{P}, \end{array}
\]

so P(F) = P.

\begin{enumerate}
\def\labelenumi{\alph{enumi})}
\setcounter{enumi}{1}
\tightlist
\item
  Let \(F \in F\) . It is clear that \(P(F) \in \mathcal{P}\) . On the other hand, F is contained in the closure of a chamber relative to \(\mathcal{H}(\text{Chap. V}, \S1, \text{no. 3, formulas (6)}),\) and so is a facet relative to the set of walls of this chamber (Chap. V, \(\S1\) , no. 4, Prop. 9). Consequently, \(\overline{F}\) is of the form \(\{v \in V \mid \alpha(v) \geq 0 \text{ for all } \alpha \in T\}\) , where T is a subset of S which we can clearly take to be equal to \(P(F)\) . Thus, \(\overline{F} = \overline{F}(P(F))\) . Q.E.D.
\end{enumerate}

If \(P \in P\) , the facet F such that \(P = P(F)\) is said to be associated to P; we denote it by \(F(P)\) . We extend these conventions to the case in which \((\mathfrak{g}, \mathfrak{h})\) is split reductive.

PROPOSITION 12. Let \(\mathcal{H}\) be the set of hyperplanes of \(\mathfrak{h}_{\mathbf{R}}\) consisting of the kernels of the roots in R. Let \(\mathcal{F}\) be the set of facets of \(\mathfrak{h}_{\mathbf{R}}\) relative to \(\mathcal{H}\) . Let \(\mathcal{S}\) be the set of parabolic subalgebras of \((\mathfrak{g},\mathfrak{h})\) . For every \(\mathfrak{p} = \mathfrak{h} + \mathfrak{g}^{\mathrm{P}}\in \mathcal{S}\) , let \(\mathrm{F}(\mathfrak{p})\) be the facet associated to P. Then \(\mathfrak{p}\mapsto \mathrm{F}(\mathfrak{p})\) is a bijection from \(\mathcal{S}\) to \(\mathcal{F}\) . If \(\mathfrak{p}_1,\mathfrak{p}_2\in \mathcal{P}\) ,

\[
\mathfrak {p} _ {1} \supset \mathfrak {p} _ {2} \Longleftrightarrow \mathrm{F} (\mathfrak {p} _ {1}) \subset \overline {{\mathrm{F} (\mathfrak {p} _ {2})}}.
\]

This follows immediately from Lemma 3.

Example. The facets corresponding to the parabolic subalgebras of \((\mathfrak{g}, \mathfrak{h})\) containing a Borel algebra b are the facets contained in the closure of the chamber associated to b (cf.~the Remark above).

PROPOSITION 13. Let \(\mathfrak{p} = \mathfrak{h} + \mathfrak{g}^{\mathrm{P}}\) be a parabolic subalgebra of \((\mathfrak{g},\mathfrak{h})\) , Q the set of \(\alpha \in \mathrm{P}\) such that \(-\alpha \notin \mathrm{P}\) , and \(\mathfrak{s} = \mathfrak{h} + \mathfrak{g}^{\mathrm{P}\cap (-\mathrm{P})}\) . Then \(\mathfrak{p} = \mathfrak{s} \oplus \mathfrak{g}^{\mathrm{Q}}\) , \(\mathfrak{s}\) is reductive in \(\mathfrak{g}\) , and \(\mathfrak{g}^{\mathrm{Q}}\) is the largest nilpotent ideal of \(\mathfrak{p}\) and the nilpotent radical of \(\mathfrak{p}\) . The centre of \(\mathfrak{p}\) is zero.

By Prop. 2, \(\mathfrak{s}\) is reductive in \(\mathfrak{g}\) and \(\mathfrak{g}^{\mathrm{Q}}\) is a nilpotent ideal of \(\mathfrak{p}\) . If \(\mathfrak{n}\) is the largest nilpotent ideal of \(\mathfrak{p}\) , \(\mathfrak{g}^{\mathrm{Q}} \subset \mathfrak{n} \subset \mathfrak{h} + \mathfrak{g}^{\mathrm{Q}}\) (Prop. 2 (i)); if \(x \in \mathfrak{n} \cap \mathfrak{h}\) , \(\operatorname{ad}_{\mathfrak{p}} x\) is nilpotent, so \(\alpha(x) = 0\) for all \(\alpha \in \mathrm{P}\) , and hence \(x = 0\) ; this proves that \(\mathfrak{n} = \mathfrak{g}^{\mathrm{Q}}\) . Since \([\mathfrak{h}, \mathfrak{g}^{\mathrm{Q}}] = \mathfrak{g}^{\mathrm{Q}}\) , the nilpotent radical of \(\mathfrak{p}\) contains \(\mathfrak{g}^{\mathrm{Q}}\) and consequently is equal to \(\mathfrak{g}^{\mathrm{Q}}\) . Let \(z = h + \sum_{\alpha \in \mathrm{P}} u_{\alpha}\) (where \(h \in \mathfrak{h}, u_{\alpha} \in \mathfrak{g}^{\alpha}\) ) be an element of the centre of \(\mathfrak{p}\) . For all \(h' \in \mathfrak{h}\) , \(0 = [h', z] = \sum \alpha(h')u_{\alpha}\) , so \(u_{\alpha} = 0\) for all \(\alpha \in \mathrm{P}\) ; it follows that \([h, \mathfrak{g}^{\beta}] = 0\) for all \(\beta \in \mathrm{P}\) , so \(h = 0\) .

